\vspace{-4mm}
\section{实习经历}
\cventry{}{\textbf{算法研究员}，\textbf{Foundation Models / Simulation}}{\textcolor{blue}{知见智能（深圳）}}{2025.9 - 至今}{}{}{建立具身数据合成与校验 Pipeline，覆盖多 Agent VQA/CoT 清洗、2D/3D 场景图推理与数据飞轮；搭建真实任务与仿真 Benchmark 结合的具身感知评测体系；基于\textbf{500万+}高质量样本使用 LoRA 和 GKD 训练/蒸馏 GSA 模型；优化 VLM 并发推理并设计可泛化 CoT 方案；评测 SOTA 具身 VLA 模型并探索改进路径。
}

\cventry{}{\textbf{算法工程师}，\textbf{Data-AML团队}}{\textcolor{blue}{字节跳动 (深圳)}}{2025.3 - 2025.9}{}{}{
负责大型模型训练平台\textbf{Forge}的运维开发与服务升级，完善日志告警、训练后同步和服务可靠性；主导\textbf{OnCallAgent}多智能体智能运维系统，集成多模态大模型、RAG 分层记忆、MCP 工具调用、模型微调、OpenAI/Ollama/vLLM 多后端、Streamlit GUI 和 Docker 部署；为小米、华为、荣耀、vivo、OPPO 等重点客户提供问题诊断与部署交付支持。}

\cventry{}{\textbf{科研实习}，\href{https://citymind.top/about-us/}{\href{http://yuxuanliang.com/}{梁宇轩教授}带领的\textbf{CityMind}}课题组}{\textcolor{blue}{香港科技大学(广州)}}{2023.9 - 2024.9}{}{}{
开展多模态学习、时空数据挖掘和时间序列预测研究；实现并优化基于图的城市交通预测模型；开发\textbf{ST-LoRA}参数高效低秩适配框架及多模态模块，用于时空预测任务。
}

\cventry{}{\textbf{科研实习}，\href{https://chenliang.tech/}{陈亮教授}带领的\textbf{GEAR}课题组}{\textcolor{blue}{中山大学}}{2023.4 - 2023.7}{}{}{
探索图学习中的长尾和OOD问题；复现和改进GNN-base模型结构和图采样方法；设计子图编码和树图变换方法。
}

\cventry{}{\textbf{研发部门算法岗}}{\textcolor{blue}{芝麻编程教育科技 (广州)}}{2021.7-2021.10}{}{}{}{信息学竞赛题目分析；制作算法竞赛培训教学资源；参与OJ平台的开发维护与评测数据生成。}
